% Clase 21 --- Los dos regímenes de la homogénea (§3.2).
% El discriminante R^2 - 4L/C decide todo: raíces reales (no oscila) o complejas
% conjugadas (oscila bajo una envolvente exponencial). Los dos casos van en la
% misma gráfica y con la misma condición inicial, que es lo que los hace
% comparables.
\begin{center}
\begin{tikzpicture}[scale=1]

  \begin{axis}[
    fisfig, width=11.0cm, height=5.2cm,
    xlabel={$t$}, ylabel={$Q_H(t)$},
    xmin=0, xmax=9, ymin=-0.75, ymax=1.15,
    xtick=\empty, ytick={0}, yticklabels={$0$},
    legend style={at={(0.98,0.95)}, anchor=north east, draw=figgray!40,
                  fill=white, font=\footnotesize},
  ]
    % envolvente
    \addplot[figgray, dashed, smooth, samples=120, domain=0:9]
      {exp(-0.42*x)};
    \addlegendentry{envolvente $\pm e^{-\gamma t}$}
    \addplot[figgray, dashed, smooth, samples=120, domain=0:9,
             forget plot] {-exp(-0.42*x)};
    % subamortiguado
    \addplot[figblue, smooth, samples=400, domain=0:9]
      {exp(-0.42*x)*cos(deg(3.1*x))};
    \addlegendentry{$R$ chica: subamortiguado}
    % sobreamortiguado
    \addplot[figred, smooth, samples=200, domain=0:9]
      {1.18*exp(-0.62*x) - 0.18*exp(-4.1*x)};
    \addlegendentry{$R$ grande: sobreamortiguado}
  \end{axis}

  \node[figlbl, anchor=north, align=center] at (5.6,-1.6)
    {$r_\pm = \dfrac{-R \pm \sqrt{R^2 - 4L/C}}{2L}$\\[0.5em]
     subamortiguado: $r_\pm = -\gamma \pm i\omega'$, \;
     $\gamma = \dfrac{R}{2L}$, \;
     $\omega' = \sqrt{\dfrac{1}{LC} - \dfrac{R^2}{4L^2}}$};

\end{tikzpicture}\\[0.5em]
{\small\itshape Toda la homogénea se decide mirando el signo del discriminante
$R^2 - 4L/C$. Si la resistencia es \textbf{grande}, las dos raíces son reales y
negativas y la carga se va a cero sin oscilar ni una vez: es el caso
\emph{sobreamortiguado}, el de una puerta con amortiguador que se cierra sola
sin golpear. Si la resistencia es \textbf{chica}, las raíces son complejas
conjugadas, $-\gamma \pm i\omega'$, y aparece lo interesante: la parte
imaginaria produce la oscilación y la real, siempre negativa, la apaga. Queda
$Q_H = Q_0 e^{-\gamma t}\cos(\omega' t + \varphi)$, es decir el $LC$ de la
primera parte de la clase \emph{modulado} por una exponencial decreciente ---la
envolvente punteada---. Dos controles que conviene hacer siempre: la frecuencia
amortiguada $\omega'$ es \emph{menor} que la natural $\omega = 1/\sqrt{LC}$, y
tiende a ella cuando $R\to 0$, donde además la envolvente se vuelve plana y se
recupera el oscilador sin pérdidas. En los dos casos $Q_H\to 0$: por eso es el
transitorio.}
\end{center}
